% NS-55 first-crossing subtask. Exact identities and scoped countermodel.
% No H^1 assumption, shape derivative, numerical window, or RH assumption.
\subsection*{First-crossing geometry on the complete form domain}

Let $V_a=\mathcal D(q_a)$ be the complete form domain from NS-51,
with unitary Fourier transform of the zero extension understood.
Let $\mu(a)=\inf\sigma(A_a)$. A \emph{first-zero window} means
\begin{equation}\label{eq:ns55fc-first}
 \mu(a_*)=0,\qquad \mu(b)>0\quad(0<b<a_*).
\end{equation}
This is a conditional setup, not an assertion that such a window exists.
If a zero or negative lower edge occurs, continuity and small-window
positivity yield such a first window. The required continuity is
Suzuki's Theorem 1.3, with its complete-domain scaling argument in
Section 4 of version 2
(\url{https://arxiv.org/html/2606.09096v2#S4}).

\begin{proposition}[Support saturation and exact inward variation]
\label{prop:ns55fc-geometry}
Assume \eqref{eq:ns55fc-first}, and let $0\ne f\in\ker A_{a_*}$.
Then the essential support of its zero extension has diameter $2a_*$;
its extreme support points are $-a_*$ and $a_*$. For $0<r<1$ set
\[
 (U_rf)(x)=r^{-1/2}f(x/r),
\]
where $f$ has already been extended by zero. Then $U_rf\in V_{ra_*}$,
$\|U_rf\|_2=\|f\|_2$, and
\begin{equation}\label{eq:ns55fc-null-variation}
 q_{ra_*}[U_rf]
 =q_{a_*}[U_rf-f]
 \ge\mu(ra_*)\|f\|_2^2>0.
\end{equation}
All identities are valid on the complete form domain. Inward dilation
preserves each parity sector.
\end{proposition}
\begin{proof}
The global compact-support form is invariant under simultaneous
translation, and its restrictions are support consistent. If the
essential support diameter were smaller than $2a_*$, a translation
would place $f$ in $V_b$ for some $b<a_*$. Its unchanged form value
zero contradicts $q_b\ge\mu(b)\|\cdot\|_2^2$.

The Fourier identity $\widehat{U_rf}(\xi)=\sqrt r\,\widehat f(r\xi)$
and comparison of $\log(2+|\xi|/r)$ with $\log(2+|\xi|)$ prove
membership in $V_{ra_*}$. Since $f$ is an operator nullvector,
$q_{a_*}(f,v)=0$ for every $v\in V_{a_*}$ by the representation
theorem for closed forms. Expand $q_{a_*}[U_rf-f]$; both cross terms
and $q_{a_*}[f]$ vanish. Support consistency and the variational
principle on $I_{ra_*}$ give \eqref{eq:ns55fc-null-variation}.
No derivative of $r\mapsto U_rf$ is taken.
\end{proof}

The endpoint conclusion means that an argument requiring an empty
exterior support strip \emph{inside} $I_{a_*}$ cannot be applied to
a first-zero nullvector without a new arithmetic step. It does not
mean that $f$ has a nonzero trace: the archived continuous zero
extension and logarithmic endpoint decay are fully compatible with
support saturation.

\begin{proposition}[Exact dilation defect with both poles and all primes]
\label{prop:ns55fc-dilation}
For $f\in V_a$, let
\[
 R_f(t)=\int_{\mathbb R}f(x+t)\overline{f(x)}\,dx,
 \quad P_f(r)=\int_{-a}^a f(x)\cosh(rx/2)\,dx,
 \quad S_f(r)=\int_{-a}^a f(x)\sinh(rx/2)\,dx,
\]
and set $\alpha(\xi)=\operatorname{Re}\psi(1/4+i\xi/2)-\log\pi$.
For $0<r<1$ define
\begin{align}
 \mathcal A_r[f]&=\int_{\mathbb R}
    [\alpha(\xi/r)-\alpha(\xi)]|\widehat f(\xi)|^2\,d\xi,
       \label{eq:ns55fc-arch}\\
 \mathcal P_r[f]&=2r\bigl(|P_f(r)|^2-|S_f(r)|^2\bigr)
                  -2\bigl(|P_f(1)|^2-|S_f(1)|^2\bigr),\notag\\
 \mathcal R_r[f]&=2\sum_{1<n<e^{2a}}\frac{\Lambda(n)}{\sqrt n}
          \operatorname{Re}\{R_f(\log n/r)-R_f(\log n)\}.
          \notag
\end{align}
Then
\begin{equation}\label{eq:ns55fc-defect}
 q_{ra}[U_rf]-q_a[f]
       =\mathcal A_r[f]+\mathcal P_r[f]-\mathcal R_r[f].
\end{equation}
For nonzero $f$,
\begin{equation}\label{eq:ns55fc-arch-bound}
 0<\mathcal A_r[f]\le5\log(1/r)\|f\|_2^2.
\end{equation}
In the even sector $S_f(r)=0$; in the odd sector $P_f(r)=0$.
Neither pole sign may be dropped in the full-space identity.
\end{proposition}
\begin{proof}
The exact complete formula, obtained from
\eqref{eq:v129-jump-decomposition} and
Proposition~\ref{prop:v135-both-parities}, is
\[
 q_a[f]=\int\alpha(\xi)|\widehat f(\xi)|^2\,d\xi
  +2|P_f(1)|^2-2|S_f(1)|^2
  -2\sum_{1<n<e^{2a}}\frac{\Lambda(n)}{\sqrt n}
                        \operatorname{Re}R_f(\log n).
\]
This extends from the smooth core by the form norm: at fixed $a$
the pole and finite prime operators are bounded, and the logarithmic
multiplier defines the principal closed form. Change variables in
the Fourier integral. Directly,
$R_{U_rf}(t)=R_f(t/r)$ and
$P_{U_rf}(1)=\sqrt r P_f(r)$, with the same formula for $S$.
The correlation $R_f$ is continuous, and $R_f(t)=0$ for
$|t|\ge2a$. Thus the dilated prime sum can use the fixed upper
cutoff $e^{2a}$: every added term past $e^{2ra}$ is exactly zero,
including equality. This proves \eqref{eq:ns55fc-defect}, retaining
every changed support overlap rather than differentiating the prime
cutoff or discarding boundary terms.

For $s_k=k+1/4$ and $t=|\xi|/2$, the convergent digamma series gives
\[
 \xi\alpha'(\xi)
 =\sum_{k\ge0}\frac{2t^2s_k}{(s_k^2+t^2)^2}>0\quad(\xi\ne0).
\]
For fixed $t>0$, $F_t(s)=2t^2s/(s^2+t^2)^2$ is nonnegative and
unimodal. Comparing its unit-spaced samples on either side of its
maximum with their adjacent integrals gives
\[
 \sum_{k\ge0}F_t(k+1/4)
 \le\int_{1/4}^\infty F_t(s)\,ds
               +2\sup_{s\ge1/4}F_t(s)\le1+4=5;
\]
the last estimate uses $F_t(s)\le1/(2s)$. Integrate this logarithmic
derivative from $|\xi|$ to $|\xi|/r$. Strict increase away from the
origin and the fact that a nonzero $L^2$ transform cannot be supported
at one point prove \eqref{eq:ns55fc-arch-bound}.
\end{proof}

At a first-zero nullvector, the exact necessary inequality is therefore
\begin{equation}\label{eq:ns55fc-necessary}
 \mathcal R_r[f]
 \le\mathcal A_r[f]+\mathcal P_r[f]
             -\mu(ra_*)\|f\|_2^2
 <\mathcal A_r[f]+\mathcal P_r[f],\qquad 0<r<1.
\end{equation}
There is no contradiction: inward scaling has positive energy.
An independent arithmetic estimate forcing the reverse non-strict
inequality for even one $r<1$ at every proposed first-zero nullvector
would exclude such a nullvector. No such estimate is derived here.
Merely naming that reverse inequality restates a sufficient exclusion
test; it does not make an arithmetic lemma out of the variational identity.

The form domain does not in general make $R_f$ differentiable at the
prime shifts. Its Fourier representation controls a logarithmic moment,
not $\int|\xi||\widehat f(\xi)|^2d\xi$. Consequently dividing
\eqref{eq:ns55fc-defect} by $1-r$ and passing to a differentiated
virial identity would require an additional regularity or limit
argument. The finite identity requires neither. Strong continuity of
dilations in the logarithmic form topology, and hence continuity of
the defect, follows by smooth approximation and locally uniform
dilation bounds in that topology; it supplies no differentiability rate.

\begin{proposition}[Translations expose a boundary interaction, not
an exclusion]\label{prop:ns55fc-translates}
Let $f\in\ker A_a$ and $h\in\mathbb R\setminus\{0\}$.
In $V_{a+|h|}$ put
$T_hf(x)=f(x-h)$ and $B_h=q_{a+|h|}(f,T_hf)$.
Both diagonal form values are zero. For every real $\theta$,
\begin{equation}\label{eq:ns55fc-translation}
 q_{a+|h|}[f\mathbin{\pm}e^{i\theta}T_hf]
          =\mathbin{\pm}2\operatorname{Re}(e^{-i\theta}B_h).
\end{equation}
Moreover, for every $v\in V_a$,
\begin{equation}\label{eq:ns55fc-boundary}
 B_h=q_{a+|h|}(f,T_hf-v).
\end{equation}
Thus a nonzero $B_h$ gives a negative direction in the larger window;
it does not contradict positivity in smaller windows. No assertion
that $B_h$ is nonzero for an arithmetic nullvector is made.
\end{proposition}
\begin{proof}
Use translation invariance and support consistency for the diagonals,
then polarize. The core null equation extends to every $v\in V_a$,
so its cross pairing vanishes and gives \eqref{eq:ns55fc-boundary}.
Choosing $\theta=\arg B_h$ makes the minus sign negative if $B_h\ne0$.
A nonzero compactly supported function and a nontrivial translate are
linearly independent, so this is a genuine nonzero test. The identity
retains the entire pairing with the portion not represented inside
$V_a$; it does not replace that interaction by an unjustified boundary
trace or omit the prime and pole terms.
\end{proof}

\begin{proposition}[A complete-domain countermodel to geometric
crossing exclusion]\label{prop:ns55fc-countermodel}
Fix a real constant $c$. On $I_a$ consider the closed restricted
logarithmic form
\[
 \ell_{a,c}[f]=\int_{\mathbb R}\log|\xi|
                     |\widehat f(\xi)|^2\,d\xi-c\|f\|_2^2.
\]
It has domain $V_a$, is support consistent and translation invariant
on compactly supported tests, and its self-adjoint operator has
compact resolvent. If $\nu_1$ is the lowest eigenvalue of
$\ell_{1,0}$, then its lowest eigenvalue is exactly
\begin{equation}\label{eq:ns55fc-model-crossing}
 \nu_c(a)=\nu_1-\log a-c.
\end{equation}
In particular it is positive for $a<e^{\nu_1-c}$, zero at that
finite first window, and negative afterwards. It has precisely
the geometric features used above and an actual first crossing.
This is an abstract operator countermodel, not the arithmetic Weil form.
\end{proposition}
\begin{proof}
This is $\tfrac12L_\Delta^{\mathrm{Dir}}-cI$, where the restricted
logarithmic Laplacian has symbol $2\log|\xi|$, as in
\cite{HLSLogBoundary}, introduction and Appendix A.1.
On a fixed interval the low-frequency negative term is bounded by
the ordinary norm: $|\widehat f(\xi)|^2\le(a/\pi)\|f\|_2^2$
and $\int_{-1}^1|\log|\xi||d\xi<\infty$. The positive logarithmic
part defines the closed form with domain $V_a$. Its embedding into
$L^2(I_a)$ is compact (Suzuki v2, Proposition 4.1), so the lower
edge is attained. The unitary dilation from $I_1$ to $I_a$ gives
\[
 \ell_{a,c}[U_af]
     =\ell_{1,0}[f]-(\log a+c)\|f\|_2^2.
\]
Take the infimum. All cross terms remain present in the associated
sesquilinear form; no finite-dimensional analogy is substituted.
\end{proof}

\paragraph{Finding and remaining input.}
First-crossing geometry provides support saturation and an exact
complete-domain dilation defect. It does not prove uniqueness. The
negative-direction calculation after a crossing describes a possible
crossing rather than excluding one, and the restricted logarithmic
countermodel makes that limitation rigorous. Any arithmetic exclusion
must use information beyond inclusion, continuity, translation, and
the logarithmic principal part. In the displayed dilation test that
information is a signed estimate for the prime-shift overlap changes
\emph{together with} the two pole changes. None is supplied by the
affine-potential equation alone in this checkpoint. Excluding kernels
only when $A_a\ge0$ would suffice for sign continuation, but no such
restricted arithmetic injectivity theorem is proved here. API, G2,
and RH remain open. No new window or tail metric was computed.
